% Section 5.3, from lectures/L05/README.md: the objectives, the questions,
% the further reading, and what the next chapter does with all of it.


\section{Review}
\label{c5:sec:review}

\needspace{6\baselineskip}
\subsection*{What you should be able to do now}
After this chapter, you should be able to:
\begin{itemize}
  \item Write down the collector current of a BJT from its base-emitter voltage, and the base current from the collector current.
  \item Say what 60 mV of extra base-emitter voltage does, and why that number is familiar.
  \item Classify an operating point as cutoff, active or saturated from two voltages.
  \item Design a switch from a load current and a logic level, without using $h_{FE}$ anywhere.
  \item Say what forced beta is, why it is used, and what it costs.
  \item Write down a MOSFET's drain current in both of its regions, and the boundary between them.
  \item Say why a MOSFET stage needs about ten times the current of a BJT stage for the same gain, and where that stops being true.
\end{itemize}

\needspace{6\baselineskip}
\subsection*{Questions to test yourself}
You should be able to answer:
\begin{enumerate}
  \item Two transistors of the same part number have $h_{FE}$ of 80 and 300. Which of your designs should care, and which should not?
  \item A transistor has 1 mA of collector current. What is the incremental resistance of its base-emitter junction, and what will that quantity be called from Chapter~\ref{c7:ch:smallsignal} onwards?
  \item Why is a saturated transistor's collector-emitter voltage not zero, and what would make it smaller?
  \item A switch is designed with a forced beta of 10. What happens if the transistor turns out to have $h_{FE}$ of 300? What if it has 40?
  \item Why does a switch waste less power than a linear stage doing the same job, and where does the power actually go in each case?
  \item A MOSFET and a BJT both carry 1 mA. Which has the higher transconductance, by how much, and what would you have to do to the MOSFET to equalise them?
  \item At what current do a BJT and a MOSFET have equal transconductance, and why is the answer suspicious?
\end{enumerate}

\needspace{6\baselineskip}
\subsection*{Further reading}
\begin{itemize}
  \item \emph{The Art of Electronics}, Horowitz and Hill, chapter 2, for the transistor introduced the same way, and chapter 3 for the MOSFET.
\end{itemize}

\needspace{6\baselineskip}
\subsection*{Looking ahead}
\begin{itemize}
  \item Chapter~\ref{c6:ch:bias} puts the device to work as an amplifier, which first means holding it still. Biasing, the quiescent point, and what an emitter resistor actually buys.
\end{itemize}
